\documentclass[12pt,a4paper]{article}

\usepackage{fontspec}
\setmainfont{Libertinus Serif}
\usepackage[greek.ancient,spanish]{babel}
\babelfont[greek]{rm}{Libertinus Serif}
\usepackage[a4paper,margin=3.2cm]{geometry}
\usepackage{microtype}
\usepackage[hidelinks]{hyperref}

% Griego politónico escrito directamente en UTF-8:
% \gr{...} para palabras sueltas; bloques con \foreignlanguage{greek}{...}
\newcommand{\gr}[1]{\foreignlanguage{greek}{#1}}
\newcommand{\stanzabreak}{\par\vspace{0.6\baselineskip}}

\setlength{\parindent}{1.25em}
\setlength{\parskip}{0.35em}
\raggedbottom


\hypersetup{
  pdftitle={Informe de lectura: Leer la Odisea en tiempos iletrados},
  pdfauthor={Claude (informe para el editor)}
}

\title{Informe de lectura\\\large\emph{Leer la Odisea en tiempos iletrados}}
\author{}
\date{Agosto de 2026}

\begin{document}

\maketitle

\section*{Ficha}

\begin{itemize}
  \item \textbf{Título:} \emph{Leer la Odisea en tiempos iletrados}
  \item \textbf{Autor:} Juan José Gómez Cadenas
  \item \textbf{Género:} ensayo literario, con memoria personal y poesía del propio autor
  \item \textbf{Extensión:} unas 70.000 palabras; introducción y diecinueve capítulos; 177 páginas en la maqueta actual
  \item \textbf{Origen:} entregas publicadas en \emph{Jot Down}, revisadas para el libro y ampliadas con capítulos nuevos (Menelao, La Gran Guerra, Eolo, Marcharse de Ítaca)
\end{itemize}

\section*{Qué es este libro}

Una lectura de la \emph{Odisea}, episodio a episodio pero sin servidumbre cronológica, trenzada con tres hilos. El primero es el comentario del poema, apoyado en la traducción en verso de Juan Manuel Rodríguez Tobal (Hiperión, 2026), que el libro cita en abundancia y enseña a leer. El segundo es la memoria del autor como lector: el adolescente que declamaba a Segalá en el Mar Menor, el chico que escuchaba el \emph{Viatge a Ítaca} de Llach en las golondrinas de Blanes, el adulto que ha tardado medio siglo en volver al texto. El tercero son los poemas del propio autor (\emph{Abandonando Ítaca}, Eolas, 2025), que proponen un ajuste de cuentas con Ulises y le dan al libro su final.

La tesis que ordena todo el material es simple y productiva: si se retiran los dioses, Ulises se queda sin coartada. Donde Homero atribuye al Olimpo, el autor devuelve la responsabilidad a los hombres, y el héroe queda a la vista como lo que el propio texto documenta: un saqueador que ataca a los cícones sin motivo, un capitán que pierde a todos sus hombres, un vengador que ordena la matanza de los pretendientes y el ahorcamiento de las esclavas. El libro no lo absuelve ni lo demoniza: lo humaniza, y en el último capítulo le ofrece una redención que Homero no contempló. La película de Nolan, estrenada en julio de 2026, funciona como percha de actualidad y piedra de toque para el contraste, no como muleta: el libro conversa con ella pero no depende de ella.

\section*{Fortalezas}

\begin{enumerate}
  \item \textbf{La voz.} Es la de un lector apasionado que escribe sin peaje académico: erudición real (griego citado y traducido, Parry, la cuestión homérica) servida con humor y con un registro coloquial deliberado, que rebaja la solemnidad justo cuando amenaza con instalarse. El «rebaño de guiris» que pace «cual vacas del sol» entre tiendas de chucherías es representativo: la ironía y el poema en la misma frase. Esta voz es el principal activo del libro y no debe planchársele ni una arruga.
  \item \textbf{El motor narrativo.} La tesis de la coartada divina no es una ocurrencia de prólogo: trabaja en cada episodio y acumula. El ataque a los cícones como pecado original, los lestrigones como «quien a hierro mata», el Hades como mala conciencia, la matanza leída junto a la boda roja de \emph{Juego de tronos}, y el capítulo sobre las esclavas ahorcadas como centro moral del libro. El lector que llega al final ha seguido un argumento, no una colección de artículos.
  \item \textbf{El uso de la traducción de Tobal.} El libro hace algo infrecuente: enseña a disfrutar una traducción en verso. Explica los compuestos (mientisagaz, beltrenza, milpensamientos), compara versiones, muestra el hexámetro en acción. Para muchos lectores será la puerta de entrada a leer la \emph{Odisea} entera, que es exactamente lo que el título promete.
  \item \textbf{El hilo autobiográfico.} No es adorno: es la respuesta a la pregunta del título. Por qué leer la \emph{Odisea} hoy se contesta mostrando lo que el poema ha hecho en una vida concreta. El último capítulo, donde las Ítacas del autor (Alba, el Mar Menor, Puerto de Sagunto) se superponen a las del héroe, cierra a la vez la memoria, la tesis y los poemas. Es el mejor capítulo del libro.
  \item \textbf{El momento.} Nolan en cartel, la traducción de Tobal recién publicada, años de reescrituras del mito en las mesas de novedades, y una base de lectores ya creada por las entregas de \emph{Jot Down}. Pocas veces un ensayo literario llega con el terreno tan abonado.
\end{enumerate}

\section*{Puntos de atención}

\begin{enumerate}
  \item \textbf{Es un híbrido, y hay que venderlo como tal.} Quien compre «una guía de la \emph{Odisea}» se encontrará además con unas memorias y un poemario; quien compre unas memorias se encontrará con filología. El libro sostiene bien la mezcla, pero la contracubierta debe anunciarla con claridad para que la sorpresa del capítulo final sea un premio y no un desvío.
  \item \textbf{Las citas largas en verso.} Hay lectores que saltan los versos por sistema, y aquí perderían la mitad del libro. Los primeros capítulos dan las herramientas para no hacerlo, y el riesgo es menor de lo que parece, pero conviene que la maqueta cuide la tipografía del verso: es materia principal, no ilustración.
\end{enumerate}

\section*{Lector}

El lector natural es el de \emph{Jot Down}: culto, curioso, sin formación clásica necesaria (todo el griego va transliterado y traducido). Funciona también para lectores jóvenes y universitarios, y no por concesión: los anclajes (el bullying del narrador adolescente, la boda roja, las peleas de gallos como poesía oral, la propia película) son parte del argumento, no cebo. Los capítulos breves y la estructura no lineal permiten además una lectura por entregas, que es como nació el texto. Lo único que el libro exige es gusto por la digresión y por la primera persona; a cambio, no exige haber leído la \emph{Odisea}: más bien deja con ganas de hacerlo.

\section*{Recomendación}

Publicar, y pronto, mientras la conversación sobre la \emph{Odisea} sigue abierta. El libro está terminado en lo sustancial, tiene una tesis que lo recorre entero, una voz reconocible de la primera a la última página y un final que justifica el conjunto. No necesita rebajarse para ser accesible: ya lo es. El trabajo editorial pendiente es de acabado, no de fondo.

\end{document}
